% Приложения. По А.3: изходен текст на програми, формати на структури от данни,
% екранни форми, получени резултати и др.
% Приложенията НЕ влизат в препоръчителния обем на изложението.
\phantomsection
\addcontentsline{toc}{section}{ПРИЛОЖЕНИЯ}
\section*{ПРИЛОЖЕНИЯ}
\label{sec:appendix}

\subsection*{Приложение А. Формат на лексикалния запис}
\addcontentsline{toc}{subsection}{Приложение А. Формат на лексикалния запис}

\TODO{описание на полетата в един лексикален запис и пример за няколко записа,
прочетени от действителния файл в data/}

\subsection*{Приложение Б. Граматика на предметната област}
\addcontentsline{toc}{subsection}{Приложение Б. Граматика на предметната област}

\TODO{пълният набор от продукции, прочетен от файла с граматиката}

\subsection*{Приложение В. Набор от въпроси за експеримента}
\addcontentsline{toc}{subsection}{Приложение В. Набор от въпроси за експеримента}

\TODO{замразеният набор от въпроси заедно с еталонните заявки за всеки от тях}

\subsection*{Приложение Г. Избрани откъси от изходния текст}
\addcontentsline{toc}{subsection}{Приложение Г. Избрани откъси от изходния текст}

Показват се само относимите откъси, а не целите файлове. Всеки откъс се чете от
истинския файл при компилация с \code{\textbackslash lstinputlisting}, тоест не може
да се разминава с хранилището.

\TODO{откъсите от src/ и include/, добавени с \textbackslash lstinputlisting след
реализацията}
